\section{Working with Git, GitHub and Codeberg}
\label{sec:git}

\subsection{Daily workflow}

The repository is a standard Git project; the default branch is
\code{main} and it should always pass both test programs. The safe cycle
for any change, however small:

\begin{verbatim}
git clone https://github.com/<owner>/SPETC.git
cd SPETC
git switch -c fix/moon-model        # a branch per change
# ... edit ...
python3 self_check.py               # both must pass before every commit
python3 tests/test_spectral_pipeline.py
git add -A
git commit -m "Fix Moon scattering normalization

One paragraph on what changed physically and why; cite the
reference if a formula changed."
git push -u origin fix/moon-model
\end{verbatim}

Then open a pull request on the web interface, from your branch into
\code{main}, and merge it after review (even self-review: reading your
own diff in the PR view catches a remarkable number of slips). Branch
names with a prefix (\code{fix/}, \code{feature/}, \code{docs/}) keep
the branch list legible. Never commit directly to \code{main}, and never
commit a red test suite: the suite is the contract that the physics
still holds.

Commit messages follow the repository's existing style: a one-line
imperative summary, a blank line, then prose paragraphs explaining the
\emph{physics or design} of the change (not a file list --- the diff
shows the files). When a commit fixes a formula, the message names the
paper.

\subsection{Keeping a fork or a second remote (Codeberg)}

To host the project on Codeberg (or any second forge) as a mirror or as
the primary, add it as a second remote --- Git has no notion of ``the''
server:

\begin{verbatim}
git remote add codeberg https://codeberg.org/<owner>/SPETC.git
git push codeberg main                    # first publication
git push codeberg --all --follow-tags     # everything, later
\end{verbatim}

To push every branch to both forges routinely, either push twice
(\code{git push origin \&\& git push codeberg}) or configure one logical
remote with two URLs:

\begin{verbatim}
git remote set-url --add --push origin \
    https://github.com/<owner>/SPETC.git
git remote set-url --add --push origin \
    https://codeberg.org/<owner>/SPETC.git
\end{verbatim}

after which a single \code{git push} updates both. Authentication on
both forges today means a personal access token used as the password
over HTTPS, or an SSH key (\code{git@github.com:...},
\code{git@codeberg.org:...}); tokens are created in the respective web
settings under \emph{Developer settings / Applications}.

\subsection{Releases}

A release is a tag plus the artefacts:

\begin{verbatim}
git tag -a v10.1 -m "SPETC v10.1"
git push origin v10.1 && git push codeberg v10.1
\end{verbatim}

Attach to the forge release page: the source archive (generated by the
forge from the tag), the three compiled documentation PDFs, and the
single-file executables of Sect.~\ref{sec:packaging} for each platform
you can build on. The release checklist: both test suites green;
\code{README.md} version references updated; documentation recompiled;
\code{etc\_user\_config.json} \emph{not} committed with your personal
values (it is runtime state --- if it shows in \code{git status},
restore it).

\subsection{What belongs in the repository}

Source, tests, documentation sources and their compiled PDFs, the
shipped \code{data/} directory, and the example/batch files. What does
not: LaTeX build intermediates (already in \code{.gitignore}),
\code{\_\_pycache\_\_}, personal configuration, large private spectra
collections (users import those locally with \code{add\_template.py}),
and one-off analysis outputs. When in doubt: if a fresh clone would not
need it to run or to rebuild the docs, leave it out.
